\documentclass{article}
\usepackage[utf8]{inputenc}
\usepackage{amsmath}
\usepackage{natbib}
\usepackage{url}
\usepackage{hyperref}
\hypersetup{colorlinks=true, citecolor=blue, linkcolor=blue, urlcolor=blue}

\begin{document}

\section{Related Work}

Estimating heterogeneous treatment effects (HTE) in survival analysis is a critical task in personalized medicine, enabling tailored treatment strategies in fields such as critical care and chronic disease management. Recent advancements have led to a variety of methodological approaches to address the challenges posed by right-censored data and complex patient heterogeneity. These approaches can be broadly categorized into four families: metalearners with censoring adjustments, tree-based non-parametric estimators, transformation-based orthogonal learning, and interpretable subgroup discovery methods. Below, we review these categories, highlighting key contributions and positioning our proposed method, Surv-CATE, within this landscape.

\subsection{Metalearners with Censoring Adjustments}

Metalearners, such as T-learners, X-learners, R-learners, and S-learners, have been adapted to handle right-censored survival data by incorporating censoring adjustments. For instance, \citet{xu2023treatment} proposed methods that utilize inverse probability of censoring weighting and doubly robust scores to mitigate the impact of censoring. These approaches allow for the application of generic regression models while accounting for incomplete observations. Additionally, \citet{zhao2015doubly} introduced doubly robust learning techniques for estimating individualized treatment effects with censored data, enhancing robustness against model misspecification. However, these methods may struggle to capture the time-varying nature of treatment effects, a limitation that Surv-CATE addresses through its generalized R-loss framework.

\subsection{Tree-Based Non-Parametric Estimators}

Tree-based methods offer flexibility in modeling complex interactions without relying on strong parametric assumptions. Causal Survival Forests, as developed by \citet{cui2023estimating}, extend the random forest framework to survival outcomes, providing robust estimation of HTE through orthogonal estimating equations. Similarly, \citet{hu2020estimating} demonstrated that nonparametric Bayesian Additive Regression Trees within the accelerated failure time model (AFT-BART-NP) achieve strong performance in terms of bias and precision. While these methods excel in handling complex data, their black-box nature limits interpretability, a critical requirement in clinical settings. In contrast, Surv-CATE employs a Mixture-of-Experts architecture to maintain interpretability while modeling heterogeneity.

\subsection{Transformation-Based Orthogonal Learning}

Transformation-based approaches aim to simplify survival analysis by mapping event times to scalar pseudo-outcomes. \citet{xu2024estimating} proposed censoring-unbiased transformations that enable the use of standard regression techniques for HTE estimation. These methods are effective in handling censoring but may discard valuable temporal information, potentially overlooking dynamic treatment effects. Surv-CATE, by contrast, retains the full time dimension, allowing for the estimation of time-varying CATEs without such reductions.

\subsection{Interpretable Subgroup Discovery}

Methods focused on interpretability seek to identify subgroups with distinct treatment effects using human-readable rules. \citet{curth2021survite} introduced SurvITE, a framework for learning HTE from time-to-event data, producing rule-based partitions that facilitate clinical decision-making. Similarly, \citet{bo2024estimating} developed a causal subgroup discovery framework that emphasizes interpretable subgroup identification in survival outcomes. These approaches are valuable but may not capture the full spectrum of heterogeneity and often rely on the proportional hazards assumption. Surv-CATE overcomes these limitations by offering a smooth, patient-level CATE surface and latent subgroup recovery without assuming proportional hazards.

\subsection{Other Machine Learning Approaches}

Beyond the above categories, machine learning techniques, including deep learning, have been explored for survival analysis, though their application to HTE estimation is less common. For example, \citet{zhu2020targeted} proposed a three-stage modular design for estimating HTE in observational survival studies, incorporating adjustments for selection and censoring biases. Deep learning methods like DeepSurv have been developed for survival prediction but are less frequently applied to HTE estimation \citep{katzman2018deepsurv}. Surv-CATE integrates machine learning principles, such as the Mixture-of-Experts architecture, to balance flexibility and interpretability, distinguishing it from purely predictive deep learning models.

\subsection{Positioning Surv-CATE}

Surv-CATE integrates elements from these diverse approaches to provide a unified framework for estimating time-varying CATEs from right-censored data. By generalizing the R-loss to survival settings, Surv-CATE avoids the proportional hazards assumption, capturing complex heterogeneity through a Mixture-of-Experts architecture. Furthermore, it employs conformal risk control to provide finite-sample prediction bands, enhancing uncertainty quantification. Compared to metalearners, Surv-CATE offers improved handling of time-varying effects; unlike tree-based methods, it prioritizes interpretability; in contrast to transformation-based approaches, it preserves temporal information; and relative to subgroup discovery methods, it provides a more granular, patient-level CATE surface. These features make Surv-CATE a versatile tool for data-driven clinical decision support, as demonstrated in its applications to critical-care trials and cardiovascular cohorts.

\bibliographystyle{plainnat}
\begin{thebibliography}{7}

\bibitem{bo2024estimating}
Bo, N., \& Ding, Y. (2024). Estimating interpretable heterogeneous treatment effect with causal subgroup discovery in survival outcomes. \emph{arXiv preprint arXiv:2409.19241}. \url{https://arxiv.org/abs/2409.19241}

\bibitem{cui2023estimating}
Cui, Y., Kosorok, M. R., Sverdrup, E., Wager, S., \& Zhu, R. (2023). Estimating heterogeneous treatment effects with right-censored data via causal survival forests. \emph{Journal of the Royal Statistical Society Series B: Statistical Methodology}, 85(2), 179--211. \url{https://doi.org/10.1093/jrsssb/qkac013}

\bibitem{curth2021survite}
Curth, A., Lee, C., \& van der Schaar, M. (2021). SurvITE: Learning heterogeneous treatment effects from time-to-event data. \emph{Advances in Neural Information Processing Systems}, 34, 26740--26753. \url{https://papers.nips.cc/paper/2021/file/e1f2d1e460f6f7d2f717fb118b4b6f60-Paper.pdf}

\bibitem{hu2020estimating}
Hu, L., Ji, J., Li

, Y., Liu, B., \& Zhang, Y. (2020). Estimating heterogeneous survival treatment effect in observational data using machine learning. \emph{arXiv preprint arXiv:2008.07044}. \url{https://arxiv.org/abs/2008.07044}

\bibitem{xu2023treatment}
Xu, Y., Ignatiadis, N., Sverdrup, E., Fleming, S., Wager, S., \& Shah, N. (2023). Treatment heterogeneity with survival outcomes. In \emph{Handbook of Matching and Weighting Adjustments for Causal Inference} (pp. 445--482). Chapman and Hall/CRC. \url{https://doi.org/10.1201/9781003102670-20}

\bibitem{xu2024estimating}
Xu, S., Cobzaru, R., Zheng, B., Finkelstein, S. N., Welsch, R. E., Ng, K., Tzoulaki, I., \& Shahn, Z. (2024). Estimating heterogeneous treatment effect from survival outcomes via (orthogonal) censoring unbiased learning. \emph{arXiv preprint arXiv:2401.11263}. \url{https://arxiv.org/abs/2401.11263}

\bibitem{zhao2015doubly}
Zhao, Y. Q., Zeng, D., Laber, E. B., Song, R., Yuan, M., \& Kosorok, M. R. (2015). Doubly robust learning for estimating individualized treatment with censored data. \emph{Biometrika}, 102(1), 151--168. \url{https://doi.org/10.1093/biomet/asu050}

\bibitem{zhu2020targeted}
Zhu, J., \& Gallego, B. (2020). Targeted estimation of heterogeneous treatment effect in observational survival analysis. \emph{Journal of Biomedical Informatics}, 107, 103474. \url{https://doi.org/10.1016/j.jbi.2020.103474}

\end{thebibliography}

\end{document}